% E.3 正方形 ABCD 边长4. P∈BC, PQ⊥AP 交 CD 于 Q. BP=x, CQ=y.
% A(0,4), B(0,0), C(4,0), D(4,4). 示例 BP=1: P(1,0). Q 在 CD: CQ=BP·(BC-BP)/AB = 1·3/4 = 0.75 → Q=(4, 0.75)
% 由相似 △ABP ~ △PCQ: y/x = CP/AB → y = x(4-x)/4
\documentclass[tikz,border=4pt]{standalone}
\usepackage{ctex}
\usepackage{amsmath}
\usetikzlibrary{calc}
\begin{document}
\begin{tikzpicture}[scale=0.8, thick]
  \coordinate (A) at (0,4);
  \coordinate (B) at (0,0);
  \coordinate (C) at (4,0);
  \coordinate (D) at (4,4);
  \coordinate (P) at (1.5,0);
  \coordinate (Q) at (4,0.9375);  % y = 1.5·2.5/4 = 0.9375

  \draw (A)--(B)--(C)--(D)--cycle;
  \draw[red, thick] (A)--(P)--(Q);

  % 直角 ∠B
  \draw ($(B)+(0.25,0)$)--($(B)+(0.25,0.25)$)--($(B)+(0,0.25)$);
  % ∠C
  \draw ($(C)+(-0.25,0)$)--($(C)+(-0.25,0.25)$)--($(C)+(0,0.25)$);
  % ∠APQ = 90°
  \draw ($(P)+(0.3,0)$) -- ($(P)+(0.3,0)+0.25*(cos(180-atan2(4,-1.5)),sin(180-atan2(4,-1.5)))$);

  \node[above left] at (A) {$A$};
  \node[below left] at (B) {$B$};
  \node[below right] at (C) {$C$};
  \node[above right] at (D) {$D$};
  \node[below] at (P) {$P$};
  \node[right] at (Q) {$Q$};

  \node[font=\footnotesize, below] at ($(B)!0.5!(P)$) {$x$};
  \node[font=\footnotesize, right] at ($(C)!0.5!(Q)$) {$y$};
\end{tikzpicture}
\end{document}
